\documentclass[12pt]{article}
\usepackage[utf8]{inputenc}
\usepackage[T1]{fontenc}
\usepackage{hyperref}
\title{Papapana Language Resources for K-12 Curriculum Development}
\author{Compiled by Claude Sonnet 5.5 (claude-sonnet-5-5) for Andr\'as Kornai}
\date{}
\begin{document}
\maketitle

\section{Introduction}
Papapana (ISO 639-3 ppn; Glottocode papa1265) is an Austronesian (Oceanic) language of Bougainville, Papua New Guinea, with roughly 120 speakers according to the sources found. It is classed as endangered. No school materials in Papapana turned up in the searches; what exists is linguistic documentation.

\section{Primary Educational Resources}
No primers, textbooks or school programmes in Papapana were found in the search results.

\section{Linguistic and Reference Materials}
\begin{itemize}
\item A Grammar of Papapana: An Oceanic Language of Bougainville, Papua New Guinea (Pacific Linguistics 659, 2020), the first comprehensive grammar: \url{https://www.harvard.com/book/9781501520730}
\item Wikipedia overview of the language: \url{https://en.wikipedia.org/wiki/Papapana_language}
\item University of Warwick project page, Describing Papapana: \url{https://warwick.ac.uk/fac/soc/al/research/describingpapapana}
\end{itemize}

\section{Religious and Community Resources}
\begin{itemize}
\item Joshua Project entry for Papapana (population and religion data only; no Papapana scripture was verified): \url{https://joshuaproject.net/people_groups/14304/PP}
\end{itemize}

\section{Digital and Audio Resources}
\begin{itemize}
\item PARADISEC catalogue, collection BP4 (recordings, reported to date from 2006): \url{https://catalog.nabu-prod.paradisec.org.au/collections/BP4}
\item ELDP archive record (documentation deposit): \url{https://eldp.access.preservica.com/dk0313/}
\end{itemize}

\section{Curriculum Development Context}
With very few speakers, mostly adults, a school curriculum would depend on revitalisation work built on the existing grammar and archived recordings. Any materials would need to be developed with the community; none are known to exist now.

\end{document}